\documentclass[11pt]{article}
\usepackage[utf8]{inputenc}
\usepackage{geometry}
\geometry{a4paper, margin=1in}
\usepackage{tabularx}
\usepackage{booktabs}
\usepackage{hyperref}
\usepackage{graphicx}
\usepackage{xcolor}

\title{\textbf{Synthetic Dataset Architecture \& Schema Documentation}\\
\large Bi-Directional Decision Support System (MCDM-ML)}
\author{Data Engineering Phase}
\date{\today}

\begin{document}

\maketitle

\section{Introduction}
This document outlines the structural schema of the generated synthetic dataset (\texttt{example\_synthetic\_data.xlsx}). The dataset acts as the foundational ``Historical Performance Data'' layer required to train Machine Learning (ML) prediction models and validate the Fuzzy Analytic Hierarchy Process (FAHP) evaluation engines. 

The workbook is divided into distinct sheets that separate static entity profiles, subjective expert weightings, longitudinal (time-series) evaluations, and operational relationship outcomes.

\section{Sheet Breakdown and Schema}

\subsection{1. Metadata Sheet}
Contains the global parameters used during the generation sequence.
\begin{itemize}
    \item \textbf{Generation Timestamp:} When the dataset was created.
    \item \textbf{Random Seed:} Ensures cryptographic reproducibility of the dataset.
    \item \textbf{Time Horizon:} The start and end years (e.g., 2018--2026) dictating the longitudinal scope of the time-series logs.
\end{itemize}

\subsection{2. Buyers \& Suppliers Sheets (Entity Profiles)}
These two sheets contain the static profiles and hidden statistical variables (Latent Factors) for every entity in the market.
\begin{table}[h]
\centering
\begin{tabularx}{\textwidth}{lX}
\toprule
\textbf{Column Name} & \textbf{Description} \\
\midrule
\texttt{entity\_id} & Unique identifier (e.g., BUY-0001, SUP-0042). \\
\texttt{industry} & The market sector (e.g., Automotive, Electronics). \\
\texttt{region} & Geographical location (used for proximity bonuses). \\
\texttt{company\_size} & Small, Medium, Large, or Enterprise (dictates capacity). \\
\texttt{creation\_year} & The year the entity entered the simulation. \\
\texttt{[latent\_factors]} & The hidden Z-scores (e.g., \textit{operational\_capability}, \textit{financial\_health}) acting as the entity's foundational DNA. \\
\bottomrule
\end{tabularx}
\end{table}

\subsection{3. AHP Weights Sheets (Buyer \& Supplier Priorities)}
These sheets store the normalized eigenvector weights representing the strategic priorities of the entities. In the real system, these are derived from FAHP pairwise matrices.
\begin{table}[h]
\centering
\begin{tabularx}{\textwidth}{lX}
\toprule
\textbf{Column Name} & \textbf{Description} \\
\midrule
\texttt{entity\_id} & The ID of the evaluating entity. \\
\texttt{[9x Criteria]} & The explicit percentage weight (summing to 1.0) assigned to each of the 9 evaluation criteria. \\
\bottomrule
\end{tabularx}
\end{table}

\subsection{4. Ratings Sheets (Time-Series Logs)}
These sheets contain the longitudinal observable scores. This is the primary \textbf{Feature Matrix (X)} for the Machine Learning model.
\begin{table}[h]
\centering
\begin{tabularx}{\textwidth}{lX}
\toprule
\textbf{Column Name} & \textbf{Description} \\
\midrule
\texttt{period\_year} & The chronological timestamp of the evaluation (e.g., 2021). \\
\texttt{entity\_id} & The company being evaluated. \\
\texttt{[9x Criteria]} & The 1.0 to 10.0 observable ratings derived from the latent factors, incorporating AR(1) temporal drift and exogenous shocks. \\
\bottomrule
\end{tabularx}
\end{table}

\subsection{5. Relationships Sheet (Operational Outcomes)}
This sheet represents the output of the Capacity-Constrained Bipartite Matching algorithm. It serves as the \textbf{Target Variable (Y)} for Machine Learning training.
\begin{table}[h]
\centering
\begin{tabularx}{\textwidth}{lX}
\toprule
\textbf{Column Name} & \textbf{Description} \\
\midrule
\texttt{relationship\_id} & Unique identifier for the contract edge. \\
\texttt{buyer\_id / supplier\_id} & The paired entities. \\
\texttt{start\_date / end\_date} & The lifecycle of the partnership. \\
\texttt{actual\_lead\_time} & Simulated physical delay (e.g., 45 days). \\
\texttt{dispute\_count} & Number of operational conflicts during the period. \\
\texttt{relationship\_quality} & Categorical label (e.g., Excellent, Weak, Terminated). \\
\texttt{collaboration\_score} & The final 1.0 to 10.0 synergy score determining success. \\
\bottomrule
\end{tabularx}
\end{table}

\section{Machine Learning Pipeline Mapping}
To fulfill the ``Machine Learning Integration'' stage of the framework:
\begin{enumerate}
    \item \textbf{Input (Features):} The ML model will ingest the \texttt{period\_year} rows from the \textbf{Ratings} sheets.
    \item \textbf{Target (Labels):} The ML model will train against the \texttt{dispute\_count} and \texttt{collaboration\_score} columns from the \textbf{Relationships} sheet.
    \item \textbf{Objective:} To predict $Y$ (Future Relationship Quality) given $X$ (Historical Time-Series Ratings).
\end{enumerate}

\end{document}
